\documentclass[12pt]{article}
\usepackage[utf8]{inputenc}
\usepackage[T1]{fontenc}
\usepackage{hyperref}
\title{Tita Language Resources for K-12 Curriculum Development}
\author{Compiled by Claude Sonnet 5 (claude-sonnet-5) for Andr\'as Kornai}
\date{}
\begin{document}
\maketitle
\section{Introduction}
Tita (ISO 639-3 tdq; also called Hoai Petel) is a language of Nigeria, spoken mainly in Taraba State (Jalingo area) and reportedly Plateau State. Wikipedia lists about 3,400 speakers (a 2000 record, with some sources giving about 6,600); Joshua Project lists a people-group population of about 7,100 and describes the language as stable. Wikipedia calls it an unclassified Benue--Congo language. The search found very little beyond these summary pages.
\section{Primary Educational Resources}
No primers, textbooks or school programmes for Tita were found in the search results. Joshua Project states that the language is not taught in schools.
\section{Linguistic and Reference Materials}
Only encyclopedic summaries were found; no grammar, dictionary or orthography was located.
\begin{itemize}
\item Wikipedia, Tita language: \url{https://en.wikipedia.org/wiki/Tita_language}
\end{itemize}
\section{Religious and Community Resources}
Joshua Project reports no Bible translation yet.
\begin{itemize}
\item Joshua Project, Tita in Nigeria: \url{https://joshuaproject.net/people_groups/15503/NI}
\item Joshua Project, language page tdq: \url{https://joshuaproject.net/index.php/languages/tdq}
\end{itemize}
\section{Digital and Audio Resources}
No audio, video, corpus or software resources were found.
\section{Curriculum Development Context}
Little documentation exists. Any curriculum would first require orthography development and community-led documentation. The English-medium Nigerian school system and neighbouring larger languages are the likely context, but this was not verified in the search.
\end{document}